\documentclass[11pt]{article}
\usepackage[a4paper,margin=1in]{geometry}
\usepackage{fontspec}
\usepackage[boldfont=false]{xeCJK}
\setCJKmainfont{FandolSong-Regular.otf}
\usepackage{amsmath}
\usepackage{amssymb}
\usepackage{array}
\usepackage{booktabs}
\usepackage{graphicx}
\usepackage{adjustbox}
\usepackage{enumitem}
\setlist[itemize]{topsep=2pt,partopsep=0pt,parsep=0pt,itemsep=2pt,leftmargin=1.4em}
\setlist[enumerate]{topsep=2pt,partopsep=0pt,parsep=0pt,itemsep=2pt,leftmargin=1.6em}
\usepackage{parskip}
\usepackage{fvextra}
\fvset{breaklines=true,breakanywhere=true,fontsize=\small}
\setlength{\parindent}{0pt}

\begin{document}

\begin{center}
  {\LARGE\bfseries Space physics}\\[0.35em]
  {\large Igcse Physics}
\end{center}
\vspace{0.6em}

\section*{The Earth, the Sun and the Moon}

The Earth spins (\textbf{rotates} 自转) on its \textbf{axis} 轴 once in about $24$ hours. The axis is \textbf{tilted} 倾斜 (not straight up). This spin gives us day and night, and makes the \textbf{Sun} 太阳 appear to move across the sky each day.

The Earth also moves around the Sun in an \textbf{orbit} 轨道, once in about $365$ days (one year). Because the axis is tilted, different parts of the Earth lean towards the Sun at different times of the year. This gives the \textbf{seasons} 季节: it is summer in the half of the Earth that leans towards the Sun.

The \textbf{Moon} 月球 orbits the Earth once in about one month. As it goes round, we see different amounts of its lit side, which gives the \textbf{phases} 月相 of the Moon (new moon, half moon, full moon).

So, from shortest to longest: one day (Earth's spin) \textless{} one month (Moon's orbit) \textless{} one year (Earth's orbit).

\section*{Orbital speed}

For an object going round a circle, the distance for one full orbit is the \textbf{circumference} 周长 $2\pi r$, where $r$ is the radius of the orbit. The \textbf{orbital speed} 轨道速率 is this distance divided by the \textbf{period} 周期 $T$ (the time for one orbit):

\[
v = \frac{2\pi r}{T}
\]

For example, a space station orbiting at $r = 7000\ \text{km}$ that takes $T = 5800\ \text{s}$ for one orbit has a speed of about $v = 2\pi(7\,000\,000)/5800 \approx 7600\ \text{m/s}$. The same equation works for planets, moons and satellites.

\section*{The Solar System}

The \textbf{Solar System} 太阳系 has one star, the Sun, at its centre. Around it move the eight \textbf{planets} 行星. In order from the Sun they are: Mercury, Venus, Earth, Mars, Jupiter, Saturn, Uranus and Neptune.

\begin{itemize}
\item The four planets nearest the Sun are small and \textbf{rocky} 岩石.

\item The four planets furthest from the Sun are large and made of \textbf{gas} 气体.

\end{itemize}

The Solar System also contains:

\begin{itemize}
\item minor planets, such as \textbf{dwarf planets} 矮行星 like Pluto, and \textbf{asteroids} 小行星 (most lie in the asteroid belt between Mars and Jupiter);

\item moons (natural \textbf{satellites} 卫星) that orbit the planets;

\item \textbf{comets} 彗星 and other small bodies.

\end{itemize}

\subsection*{Orbits and gravity}

The Sun holds most of the mass of the Solar System. Its \textbf{gravity} 引力 (gravitational pull) reaches out and keeps the planets in their orbits. This force is the gravitational attraction of the Sun.

The \textbf{gravitational field strength} 重力场强度 tells you how strong gravity is:

\begin{itemize}
\item at the surface of a planet it is bigger for a planet with more mass;

\item around a planet it gets weaker as you move further away.

\end{itemize}

In the same way, the Sun's gravity gets weaker further out, so the outer planets move more slowly (smaller orbital speed) than the inner ones.

Most orbits are not perfect circles but \textbf{ellipses} 椭圆 (a stretched circle), and the Sun is not at the centre of the ellipse. A comet has a very stretched orbit. An object in an elliptical orbit moves \textbf{faster} when it is closer to the Sun. We explain this with the \textbf{conservation of energy} 能量守恒 — the total \textbf{energy} 能量 stays the same, so as the object's \textbf{gravitational potential energy} 重力势能 falls (coming closer), its \textbf{kinetic energy} 动能 rises (it speeds up).

\subsection*{Light travel time}

Distances in space are huge, so we often work out how long light takes to cross them. Light travels at $3.0 \times 10^8\ \text{m/s}$. Using $\text{time} = \dfrac{\text{distance}}{\text{speed}}$, light from the Sun (about $1.5 \times 10^{11}\ \text{m}$ away) takes about $500\ \text{s}$, which is roughly $8$ minutes, to reach the Earth.

\section*{How the Solar System formed}

The planets formed from a giant cloud of \textbf{gas} and \textbf{dust} 尘埃 in space, which contained many chemical \textbf{elements} 元素. This is the \textbf{accretion} 吸积 model:

\begin{itemize}
\item gravity pulled the cloud together;

\item the cloud was spinning, so it flattened into a spinning \textbf{accretion disc} 吸积盘;

\item most of the matter fell to the centre and became the Sun, while the planets grew from the leftover material in the disc.

\end{itemize}

Near the hot young Sun, only rock and metal could stay solid, so the inner planets are rocky. Far out, where it was cold, gases could also collect, so the outer planets grew large and gaseous.

\section*{The Sun as a star}

The Sun is a \textbf{star} 恒星 of medium size, made mostly of \textbf{hydrogen} 氢 and \textbf{helium} 氦. It radiates most of its \textbf{energy} in the infrared, visible light and \textbf{ultraviolet} 紫外线 parts of the \textbf{electromagnetic spectrum} 电磁波谱.

A star is powered by \textbf{nuclear reactions} 核反应 in its core. In a stable star, the reaction is \textbf{nuclear fusion} 核聚变: hydrogen nuclei join to make helium, releasing huge amounts of energy.

\section*{Stars, galaxies and the Universe}

A \textbf{galaxy} 星系 is a group of many billions of stars held together by gravity. Our Sun is one star in the galaxy called the \textbf{Milky Way} 银河系. All the other stars in the Milky Way are very much further from the Earth than the Sun is.

Distances between stars are so big that we measure them in \textbf{light-years} 光年. One light-year is the distance light travels in one year (in the vacuum of space), which is about $9.5 \times 10^{15}\ \text{m}$.

The Milky Way is about $100\,000$ light-years across (its \textbf{diameter} 直径). It is just one of many billions of galaxies. All of these galaxies together make up the \textbf{Universe} 宇宙.

\section*{The life cycle of a star}

A star is born, lives and dies over a very long time:

\begin{enumerate}
\item A star forms from an \textbf{interstellar cloud} 星际云 of gas and dust that contains hydrogen.

\item Gravity pulls the cloud inwards, and it heats up. This collapsing, heating cloud is a \textbf{protostar} 原恒星.

\item The protostar becomes a \textbf{stable} 稳定 star when the inward pull of gravity is balanced by an outward push caused by the very high temperature in its centre.

\item In time, every star runs out of hydrogen \textbf{fuel} 燃料 for fusion.

\end{enumerate}

What happens next depends on the mass of the star:

\begin{itemize}
\item A medium star (like the Sun) swells into a \textbf{red giant} 红巨星. It then throws off its outer layers as a glowing cloud called a planetary \textbf{nebula} 星云, leaving a small, hot \textbf{white dwarf} 白矮星 at the centre.

\item A much heavier star swells into a \textbf{red supergiant} 红超巨星 and then explodes as a \textbf{supernova} 超新星. This blast spreads out a nebula containing hydrogen and new, heavier elements, and leaves behind a \textbf{neutron star} 中子星 or a \textbf{black hole} 黑洞.

\end{itemize}

The nebula from a supernova can later form new stars, with planets orbiting them — so our own Solar System came from earlier stars.

\section*{The expanding Universe}

When a star or galaxy moves away from us (it is \textbf{receding} 退行), the light we receive from it has a longer \textbf{wavelength} 波长 than normal — it is shifted towards the red end of the spectrum. This is \textbf{redshift} 红移.

The light from distant galaxies is redshifted, and the further away a galaxy is, the bigger its redshift. This tells us that the galaxies are moving apart: the Universe is \textbf{expanding} 膨胀. Running this backwards, everything was once together at a single point — this is the evidence for the \textbf{Big Bang} 大爆炸 theory.

\subsection*{More evidence and the age of the Universe}

Faint \textbf{microwave} 微波 \textbf{radiation} 辐射 is found coming from every direction in space. This is the \textbf{cosmic microwave background radiation} 宇宙微波背景辐射 (CMBR). It was made as high-energy radiation soon after the Big Bang, and has been stretched out into the microwave part of the spectrum as the Universe expanded.

For a distant galaxy:

\begin{itemize}
\item its speed $v$ of moving away can be found from the redshift of its starlight;

\item its distance $d$ can be found from the brightness of a \textbf{supernova} seen in it.

\end{itemize}

The \textbf{Hubble constant} 哈勃常数 $H_0$ links these two:

\[
H_0 = \frac{v}{d}
\]

Its value today is about $2.2 \times 10^{-18}$ per second. Turning this around gives an estimate for the age of the Universe:

\[
\frac{d}{v} = \frac{1}{H_0}
\]

This works because every galaxy seems to have started from the same single point at the same time — more support for the Big Bang.


\end{document}
